% Chapter 3: Materials and Method
\chapter{Materials and Method}
\label{ch:mutbench}

This chapter presents the methodology of MutBench, an experimental framework for analyzing which factors---biological information type, detection method category, and pathogen characteristics---influence viral mutation hotspot detection performance.
The overall pipeline processes multiple sequence alignment (MSA) data from 11 RNA viruses through 20 scoring types and 14 detection method families spanning 5 categories (threshold-based, signal processing, density/clustering, statistical, and adaptive window), with 39 parameter variants, evaluating results against a 3-layer ground truth system using the hotspot-score composite metric and Matthews Correlation Coefficient (MCC).
The evaluation follows a two-stage main design with a Stage~3 information-integration layer: Stage~1 conducts in-depth analysis on SARS-CoV-2 to validate the framework components, Stage~2 extends the evaluation to all 11 pathogens for the large-scale cross-pathogen MCC comparison, and Stage~3 (Section~\ref{subsec:stage3_methods}) decomposes the resulting performance by information type via per-feature AUC, ablation, Random Forest, and vaccine-escape cross-validation.
Section~\ref{sec:mutbench_materials} describes the materials (target pathogens and sequence data).
Section~\ref{sec:mutbench_methods} presents the methods, beginning with a framework overview (Section~\ref{subsec:framework_overview}) that describes the overall architecture and three-component pipeline, followed by the 3-layer ground truth framework (Section~\ref{subsec:3layer_gt}).
The experimental design is then presented in two stages: Stage~1 (Section~\ref{subsec:stage1}) focuses on SARS-CoV-2 in-depth analysis including MutClust variant comparison and the hotspot-score metric, while Stage~2 (Section~\ref{sec:stage2}) extends to the 11-pathogen cross-pathogen benchmark with 20 scoring types and 14 detection families.
Section~\ref{sec:eval_stat_design} details the evaluation metrics and statistical analysis design.
\label{sec:9pathogen_benchmark}

%% ============================================================
\section{Materials}
\label{sec:mutbench_materials}

\subsection{Target Pathogens and Data}
\label{subsec:9pathogen_data}

Eleven RNA viruses were selected to cover diverse evolutionary dynamics, transmission routes, and genomic characteristics:
(1)~\textit{Respiratory viruses}: SARS-CoV-2 (Spike, 1,259 AA), MERS-CoV (Spike, 1,330 AA), RSV (F protein, 550 AA), Influenza A H3N2 (HA, 543 AA), and Influenza B (HA, 560 AA);
(2)~\textit{Blood-borne/sexually transmitted}: HIV-1 (Env gp120, 450 AA) and HCV (E2, 340 AA);
(3)~\textit{Enteric}: Norovirus (VP1, 536 AA);
(4)~\textit{Arthropod-borne}: Dengue (E protein, 490 AA);
(5)~\textit{Neurotropic}: Rabies (G protein, 491 AA) and EV-A71 (VP1, 261 AA).
\textit{Note: all amino acid lengths above refer to post-truncation alignment lengths used in the benchmark; full-length protein sequences may differ (e.g., SARS-CoV-2 Spike full length is 1{,}273 AA, MERS-CoV Spike 1{,}353 AA, Rabies G 524 AA, EV-A71 VP1 297 AA).}
These 11 pathogens span substitution rates from $\sim$10$^{-3}$ (HIV-1, SARS-CoV-2) to $\sim$10$^{-4}$ (Rabies) substitutions/site/year, alignment lengths from 261 to 1,330 AA, and evolutionary patterns from rapid antigenic drift (H3N2) to epochal evolution (Norovirus).
This diversity ensures that the benchmark evaluates method performance across a representative range of viral evolutionary regimes rather than a single pathogen's characteristics.
Zika virus was considered but excluded from the Stage~2 benchmark because only 52 unique E protein sequences were available in GenBank, insufficient for reliable per-position frequency and entropy estimation; Zika is included in the supplementary 12-pathogen feature analysis (Section~\ref{sec:information_analysis}).

\medskip
\noindent\textbf{Per-pathogen descriptions.}
The following paragraphs describe each pathogen in the order listed above, grouped by transmission route.

\medskip
\noindent\textit{Respiratory viruses.}

\paragraph{SARS-CoV-2 (Spike, 1,259 AA).}
A betacoronavirus (family \textit{Coronaviridae}) that caused the COVID-19 pandemic (2019--present), SARS-CoV-2 serves as the primary reference pathogen for this benchmark.
The Spike protein mediates ACE2 receptor binding and membrane fusion for cell entry, and is the principal target of vaccines and neutralizing antibodies, making it the most extensively characterized viral surface protein in history.
SARS-CoV-2 evolves at ${\sim}1 \times 10^{-3}$ substitutions/site/year for the Spike gene, with its relatively shallow pandemic phylogeny ($\sim$5 years in our dataset, 2019--2024) creating strong founder effects (e.g., D614G) while convergent evolution at immune evasion sites (E484K, N501Y) provides clear Layer~A ground truth.
The low unique sequence ratio (753/4,982 = 15.1\%) reflects massive sequencing of identical or near-identical lineages during the early pandemic.
DMS data are available from Dadonaite et al.\ (2024), covering the full Spike protein entry assay.
SARS-CoV-2 was selected as the anchor pathogen because it offers the richest ground truth data across all three layers.

\paragraph{MERS-CoV (Spike, 1,330 AA).}
A betacoronavirus (family \textit{Coronaviridae}) first identified in Saudi Arabia in 2012 and causing Middle East Respiratory Syndrome, MERS-CoV has a case fatality rate of $\sim$35\%, far exceeding that of SARS-CoV-2.
The Spike protein binds the DPP4 (CD26) receptor for cell entry; its receptor-binding domain is the primary target for therapeutic antibody development.
MERS-CoV evolves more slowly than SARS-CoV-2, with a substitution rate of ${\sim}0.9 \times 10^{-3}$ substitutions/site/year, and sporadic zoonotic spillovers from dromedary camels produce limited human-to-human transmission chains.
The moderate unique ratio (530/1,311 = 40.4\%) reflects both camel reservoir diversity and repeated independent introductions.
MERS-CoV was selected to provide a coronavirus comparator with a different receptor target and a fundamentally different epidemiological pattern (sporadic zoonotic vs.\ sustained pandemic).

\paragraph{RSV (F protein, 550 AA).}
Respiratory Syncytial Virus (family \textit{Pneumoviridae}) is the leading cause of lower respiratory tract infections in infants and a major cause of hospitalization in the elderly.
The Fusion (F) protein mediates membrane fusion during cell entry and is the target of the recently approved monoclonal antibody nirsevimab and the prefusion-stabilized vaccines (2023).
RSV-A and RSV-B co-circulate with a nucleotide substitution rate of ${\sim}1.5 \times 10^{-3}$ substitutions/site/year in the F gene; antigenic variation is more limited than influenza, with most escape mutations concentrated in antigenic site~0.
The high unique ratio (4,779/5,070 = 94.3\%) indicates substantial sequence diversity across global surveillance, with minimal identical sequences despite the relatively conserved F protein.
RSV was selected because its high sequence count and recent vaccine development provide both statistical power and clinical relevance for benchmarking.

\paragraph{Influenza A H3N2 (HA, 543 AA).}
Influenza A H3N2 (family \textit{Orthomyxoviridae}) has circulated in humans since the 1968 pandemic and remains a major driver of seasonal influenza morbidity and mortality.
The hemagglutinin (HA) protein mediates sialic acid receptor binding and is the primary target of neutralizing antibodies and seasonal vaccines; its five antigenic sites (A--E) accumulate substitutions that necessitate biannual vaccine strain updates.
H3N2 exhibits rapid antigenic drift at ${\sim}4 \times 10^{-3}$ substitutions/site/year in HA, one of the highest rates among RNA viruses, driven by strong immune selection on the globular head domain.
The very high unique ratio (2,562/2,644 = 96.9\%) reflects decades of antigenic diversification with few identical sequences.
H3N2 was selected as the archetypal rapidly drifting pathogen whose well-characterized antigenic sites provide robust Layer~A ground truth.

\paragraph{Influenza B (HA, 560 AA).}
Influenza B virus (family \textit{Orthomyxoviridae}) co-circulates with influenza A and contributes substantially to seasonal epidemic burden, particularly in children.
The HA protein serves the same receptor-binding and membrane fusion roles as in influenza A, but Influenza B evolves more slowly (${\sim}2 \times 10^{-3}$ substitutions/site/year) with two co-circulating lineages (Victoria and Yamagata) that diverged in the 1980s.
The very high unique ratio (5,019/5,325 = 94.3\%) reflects long-term co-circulation of both lineages and accumulation of distinct mutations.
Influenza B was selected to contrast with H3N2: both are segmented negative-sense RNA viruses targeting HA, but Influenza B's slower drift rate and dual-lineage structure test whether detection methods generalize across different evolutionary tempos within the same virus family.

\medskip
\noindent\textit{Blood-borne/sexually transmitted viruses.}

\paragraph{HIV-1 (Env gp120, 450 AA).}
Human Immunodeficiency Virus type~1 (family \textit{Retroviridae}) is the causative agent of AIDS and remains a major global health burden with $\sim$39 million people living with HIV.
The Env gp120 subunit mediates CD4 receptor and co-receptor (CCR5/CXCR4) binding; its extraordinary sequence variability, driven by an error-prone reverse transcriptase, has prevented effective vaccine development for decades.
HIV-1 mutates at ${\sim}2.5 \times 10^{-3}$ substitutions/site/year (one of the highest rates among the benchmark pathogens; EV-A71 is comparable~\cite{sanjuan2010viral}), with extensive within-host and between-host diversity generating up to 35\% amino acid divergence across subtypes.
The unique ratio (2,256/3,000 = 75.2\%) reflects high inter-subtype diversity.
HIV-1 was selected as the extreme end of the mutation rate spectrum, testing whether detection methods remain effective under very high sequence diversity where nearly every position is variable.

\paragraph{HCV (E2, 340 AA).}
Hepatitis C Virus (family \textit{Flaviviridae}) causes chronic liver infection affecting $\sim$58 million people worldwide, leading to cirrhosis and hepatocellular carcinoma.
The E2 glycoprotein mediates CD81 receptor binding and is the primary target of neutralizing antibodies; its hypervariable region~1 (HVR1, residues 384--410) is the most variable segment in the entire HCV genome and serves as a decoy for the immune system.
HCV evolves at ${\sim}1 \times 10^{-3}$ substitutions/site/year, with diversifying selection concentrated in HVR1 while the CD81 binding residues remain conserved.
The moderate unique ratio (1,067/1,362 = 78.3\%) reflects inter-genotype diversity across the seven major HCV genotypes.
HCV was selected because its well-defined hypervariable and conserved regions provide clear positive and negative controls for hotspot detection accuracy.

\medskip
\noindent\textit{Enteric virus.}

\paragraph{Norovirus (VP1, 536 AA).}
Norovirus (family \textit{Caliciviridae}) is the leading cause of acute viral gastroenteritis worldwide, responsible for $\sim$685 million cases annually.
The VP1 capsid protein forms the viral shell and contains the protruding (P) domain that binds histo-blood group antigens (HBGAs) for cell attachment; antigenic variation in the P2 subdomain drives immune escape.
Norovirus GII.4 exhibits epochal evolution, with long periods of stasis punctuated by the emergence of antigenically novel pandemic variants every 2--4 years, at a substitution rate of ${\sim}2 \times 10^{-3}$ substitutions/site/year.
The moderate unique ratio (791/1,500 = 52.7\%) reflects this epochal pattern, where many sequences within each epoch are nearly identical.
Norovirus was selected to test detection method performance under epochal evolutionary dynamics, a qualitatively different regime from the continuous antigenic drift of influenza.

\medskip
\noindent\textit{Arthropod-borne virus.}

\paragraph{Dengue (E protein, 490 AA).}
Dengue virus (family \textit{Flaviviridae}) is transmitted by \textit{Aedes} mosquitoes and causes $\sim$390 million infections per year across tropical and subtropical regions, with severe forms including hemorrhagic fever and shock syndrome.
The E protein mediates receptor binding and membrane fusion; it contains three structural domains (DI--DIII) with domain~III harboring the primary neutralizing epitopes that differ across the four serotypes (DENV-1 through DENV-4).
Dengue evolves at ${\sim}7.6 \times 10^{-4}$ substitutions/site/year, slower than the respiratory viruses, with inter-serotype divergence reaching $\sim$30\% at the amino acid level while intra-serotype diversity is more limited.
The moderate unique ratio (796/2,000 = 39.8\%) reflects the inclusion of multiple serotypes with substantial inter-serotype but limited intra-serotype diversity.
Dengue was selected because the four-serotype structure provides a natural test of whether detection methods can distinguish serotype-defining positions from true adaptive hotspots.

\medskip
\noindent\textit{Neurotropic viruses.}

\paragraph{Rabies (G protein, 491 AA).}
Rabies virus (family \textit{Rhabdoviridae}) causes fatal encephalitis in mammals with a near-100\% case fatality rate once symptoms appear, killing $\sim$59,000 people annually, predominantly in Asia and Africa.
The G protein is the sole surface glycoprotein, mediating receptor binding (nicotinic acetylcholine receptor, NCAM, p75NTR) and membrane fusion; it is the only target of virus-neutralizing antibodies and the sole component of current vaccines.
Rabies evolves slowly at ${\sim}4 \times 10^{-4}$ substitutions/site/year, reflecting strong purifying selection on the G protein and limited immune-driven diversification due to the virus's rapid lethality (hosts die before mounting sustained immune responses).
The unique ratio is 100\% (2,038/2,038), indicating that every sampled sequence is distinct, consistent with long-term diversification across geographically isolated reservoir hosts (dogs, bats, raccoons).
Rabies was selected as a slowly evolving pathogen with strong purifying selection, testing whether detection methods can identify the few positively selected sites against a highly conserved background.
DMS data are available for the G protein.

\paragraph{EV-A71 (VP1, 261 AA).}
Enterovirus A71 (family \textit{Picornaviridae}) causes hand, foot, and mouth disease in children, with periodic outbreaks of severe neurological complications including brainstem encephalitis and acute flaccid paralysis, particularly in the Asia-Pacific region.
The VP1 capsid protein contains the major neutralization epitopes including the BC loop, and determines viral serotype and receptor binding specificity.
EV-A71 evolves at ${\sim}4.5 \times 10^{-3}$ substitutions/site/year, comparable to influenza, with frequent genotype replacement (B1--B5, C1--C5) driven by population immunity.
The unique ratio is 100\% (662/662) and the protein is the shortest in the benchmark at 261~AA (post-truncation), providing a test case for detection methods on short, highly diverse sequences where statistical power is limited.
EV-A71 was selected to evaluate detection performance at the lower end of alignment length, where the small number of positions amplifies the difficulty of distinguishing true hotspots from noise.
DMS data are available for VP1.

\medskip
FASTA sequences for each pathogen were collected from NCBI GenBank by querying each target protein name (e.g., ``SARS-CoV-2 spike glycoprotein'') with filters for complete coding sequences, excluding partial sequences and laboratory-derived constructs. No date range restriction was applied; all available sequences as of December 2024 were downloaded. Sequences with $>$5\% ambiguous bases (N) were removed, and identical sequences were deduplicated (the ``Unique seq.'' column in Table~\ref{tab:pathogen_data} reflects post-deduplication counts). Multiple sequence alignment (MSA) was performed for each surface protein (SARS-CoV-2: Spike, H3N2: HA, HIV-1: Env gp120, etc.) using MAFFT v7.505~\cite{katoh2013mafft} with the \texttt{--auto} mode, followed by computation of per-position mutation frequency and entropy.
The \texttt{--auto} mode automatically selects among FFT-NS-2, L-INS-i, and G-INS-i based on the number and length of input sequences~\cite{katoh2013mafft}.
At the sequence scales used in this study (1,311--5,325 sequences, Table~\ref{tab:pathogen_data}), FFT-NS-2 is typically selected, but the alignment algorithm may differ across pathogens, and the effect of this difference on alignment quality---and consequently on per-position mutation frequency computation---was not controlled in this study.
When alignment lengths differed across sequences, all sequences were truncated to the minimum sequence length within the MSA to ensure consistency in per-position comparisons.
Ground truth positions exceeding the alignment length are thereby excluded from evaluation (the position counts in Table~\ref{tab:gt_positions} are effective values after this truncation).
Table~\ref{tab:pathogen_data} summarizes the characteristics of each pathogen's sequence dataset.

\begin{table}[htbp]
\centering
\caption{Sequence dataset characteristics for 11 pathogens used in the benchmark}
\label{tab:pathogen_data}
\footnotesize
\adjustbox{max width=\textwidth}{
\begin{tabular}{lllrrr}
\toprule
\textbf{Pathogen} & \textbf{Surface protein} & \textbf{Reference accession} & \textbf{Sequences} & \textbf{Unique seq.} & \textbf{Align. len. (AA)} \\
\midrule
SARS-CoV-2 & Spike & NC\_045512.2 (Wuhan-Hu-1) & 4,982 & 753 & 1,259 \\
H3N2 & HA & CY045572 (A/Brisbane/10/2007) & 2,644 & 2,562 & 543 \\
Norovirus & VP1 & KC576909 (GII.4 Sydney) & 1,500 & 791 & 536 \\
HIV-1 & Env gp120 & K03455 (HXB2) & 3,000 & 2,256 & 450 \\
Dengue & E protein & NC\_001474.2 (DENV-2) & 2,000 & 796 & 490 \\
RSV & F protein & KT992094 (RSV-A) & 5,070 & 4,779 & 550 \\
Influenza B & HA & CY018765 (B/Florida/4/2006) & 5,325 & 5,019 & 560 \\
MERS & Spike & NC\_019843.3 (MERS-CoV) & 1,311 & 530 & 1,330 \\
HCV & E2 & NC\_004102.1 (HCV-1a) & 1,362 & 1,067 & 340 \\
Rabies & G protein & NC\_001542.1 & 2,038 & 2,038 & 491 \\
EV-A71 & VP1 & AB204852 & 662 & 662 & 261 \\
\bottomrule
\end{tabular}
}
\end{table}

%% ============================================================
\section{Methods}
\label{sec:mutbench_methods}

\subsection{Framework Overview}
\label{subsec:framework_overview}

The MutBench framework follows a three-component pipeline architecture (Figure~\ref{fig:mutbench_pipeline}): (1)~ground truth construction from three independent biological evidence sources, (2)~systematic scoring and detection via 20 scoring types and 14 detection method families, and (3)~multi-metric evaluation with factorial statistical analysis.
The evaluation is organized into two stages: Stage~1 performs in-depth analysis on SARS-CoV-2 to validate the framework components, and Stage~2 extends the evaluation to all 11 pathogens for large-scale cross-pathogen comparison.
The following subsections describe each component in detail.

\begin{figure}[htbp]
\centering
\resizebox{\textwidth}{!}{%
\begin{tikzpicture}[
    node distance=1.2cm,
    box/.style={rectangle, draw, rounded corners, minimum width=3.5cm, minimum height=0.8cm, align=center, font=\small},
    arrow/.style={->, thick, >=stealth},
    label/.style={font=\footnotesize\itshape, text=gray}
]
% Input
\node[box, fill=blue!10] (input) {MSA Sequences\\(11 pathogens)};

% Scoring
\node[box, fill=green!10, below=of input] (scoring) {Scoring Types\\(6 categories, 20 types)};

% Detection
\node[box, fill=orange!10, below=of scoring] (detection) {Detection Methods\\(39 methods, 14 families)};

% Ground Truth
\node[box, fill=red!10, right=2.5cm of scoring] (gt) {3-Layer Ground Truth\\A: Convergent / B: Constrained / C: DMS};

% Evaluation
\node[box, fill=purple!10, below=of detection] (eval) {Evaluation\\(MCC, hotspot-score)};

% Statistical Analysis
\node[box, fill=yellow!10, below=of eval] (stats) {Statistical Analysis\\(ANOVA, Friedman, LOPO)};

% Arrows
\draw[arrow] (input) -- (scoring);
\draw[arrow] (scoring) -- (detection);
\draw[arrow] (detection) -- (eval);
\draw[arrow] (gt) -- (eval);
\draw[arrow] (eval) -- (stats);

% Stage labels
\node[label, left=0.3cm of scoring] {Stage 1: SARS-CoV-2};
\node[label, left=0.3cm of stats] {Stage 2: 11 pathogens};

\end{tikzpicture}
}% end resizebox
\caption{Overview of the MutBench experimental framework pipeline. Input multiple sequence alignment (MSA) data are processed through scoring types and detection methods, evaluated against a 3-layer ground truth, and analyzed using multiple statistical tests.}
\label{fig:mutbench_pipeline}
\end{figure}

\textbf{Component 1: Ground Truth Construction.}
A 3-layer ground truth system captures three orthogonal dimensions of mutational importance: Layer~A (literature-curated adaptive/diversifying positions) defines true positives, Layer~B (purifying selection) defines constrained negatives for a separate false-positive audit, and Layer~C (DMS fitness) provides experimental validation.
For Stage~1, synthetic data with known signal positions and DMS-based functional site annotations provide additional controlled validation.
Details are presented in Section~\ref{subsec:3layer_gt}.

\textbf{Component 2: Scoring and Detection.}
Twenty scoring types organized into six categories (frequency-based, MSA-derived, phylogenetic, structural, AI-based, composite) assign per-position scores from diverse biological information sources.
These scores are then processed by 14 detection method families (39 parameter variants) spanning five algorithmic categories (threshold-based, signal processing, density/clustering, statistical, adaptive window) to convert continuous score profiles into discrete hotspot calls.
Stage~1 provides a controlled SARS-CoV-2 sanity check for MutClust variants and temporal scoring (Section~\ref{subsec:stage1}); Stage~2 extends to the full scoring $\times$ detection $\times$ pathogen factorial design (Section~\ref{sec:stage2}).

\textbf{Component 3: Evaluation and Statistical Analysis.}
Position-level evaluation metrics (MCC, F1, enrichment ratio, constrained FPR, DMS F1) quantify detection performance, and factorial statistical methods (three-way ANOVA with $\omega^2$, Friedman test, LOPO cross-validation) decompose the sources of performance variation.
Stage~1 uses the hotspot-score composite metric; Stage~2 adopts MCC as the primary metric.
Details are presented in Section~\ref{sec:eval_stat_design}.

%% ============================================================
\subsection{3-Layer Ground Truth Framework}
\label{subsec:3layer_gt}

To mitigate single ground truth bias, a 3-layer ground truth framework based on three independent sources of biological evidence was constructed.
Figure~\ref{fig:3layer_gt_concept} illustrates the structure and integration of the three layers.

\begin{figure}[htbp]
\centering
\resizebox{\textwidth}{!}{%
\begin{tikzpicture}[
    node distance=1.2cm and 1.8cm,
    layer/.style={rectangle, draw, rounded corners, minimum width=4.5cm, minimum height=1.2cm, align=center, font=\small, thick},
    combine/.style={rectangle, draw, rounded corners, minimum width=5.5cm, minimum height=1.0cm, align=center, font=\small, thick},
    role/.style={font=\footnotesize, align=center},
    arrow/.style={->, thick, >=stealth}
]

% Layer A
\node[layer, fill=blue!15] (layerA) {\textbf{Layer A: Adaptive Positives}\\Literature-curated sites\\(True Positives)};

% Layer B
\node[layer, fill=green!15, right=2.5cm of layerA] (layerB) {\textbf{Layer B: Purifying Selection}\\Constrained functional sites\\(True Negatives)};

% Layer C
\node[layer, fill=orange!15, right=2.5cm of layerB] (layerC) {\textbf{Layer C: DMS Fitness}\\Experimental functional impact\\(Independent validation)};

% Role labels
\node[role, below=0.2cm of layerA] (roleA) {Evaluates recall\\(detection sensitivity)};
\node[role, below=0.2cm of layerB] (roleB) {Evaluates specificity\\(constrained FPR)};
\node[role, below=0.2cm of layerC] (roleC) {Evaluates biological\\relevance (DMS F1)};

% Integration node
\node[combine, fill=purple!10, below=2.2cm of layerB] (integrate) {\textbf{Multi-Layer Evaluation}\\MCC (Layer A) + Constrained FPR (Layer B) + DMS F1 (Layer C)};

% Arrows
\draw[arrow, blue!70] (layerA.south) -- ++(0,-1.2) -| (integrate.north west);
\draw[arrow, green!60!black] (layerB.south) -- (integrate.north);
\draw[arrow, orange!70!black] (layerC.south) -- ++(0,-1.2) -| (integrate.north east);

\end{tikzpicture}
}% end resizebox
\caption{3-Layer Ground Truth Framework. Layer~A (literature-curated adaptive/diversifying positions) defines true positives for MCC computation, Layer~B (purifying selection) defines constrained negatives for constrained FPR, and Layer~C (DMS fitness) provides experimental validation. The three layers capture orthogonal dimensions of mutational importance.}
\label{fig:3layer_gt_concept}
\end{figure}

\textbf{Design rationale: three orthogonal dimensions.}
The three layers are deliberately designed to capture different dimensions of mutational importance.
Layer~A asks whether a position is supported by literature as adaptively or diversifyingly important for that pathogen---through recurrent emergence, immune-escape evidence, epochal-variant membership, or hypervariable-region membership depending on the pathogen.
Layer~B asks ``is this position structurally constrained?'' (a \textit{structural} question), identifying positions where any mutation destroys protein function.
Layer~C asks ``what is the measurable functional impact of mutations at this position?'' (a \textit{functional} question), using laboratory DMS experiments that directly quantify how each amino acid substitution affects viral fitness (e.g., cell entry efficiency, receptor binding affinity).
This separation means each layer provides a non-redundant perspective: Layer~A captures literature-curated adaptive or diversifying evidence that DMS cannot always provide, while Layer~C captures functional effects that may not yet be visible in surveillance-derived annotations.

\subsubsection{Layer A: Adaptive Positive Selection}

Layer~A positions are literature-curated adaptive or diversifying positions treated as hotspot positives.
For pathogens where convergent-evolution evidence is available, positions where mutations recurrently arose in multiple independent lineages \textit{within the same viral species} are classified as hotspots under positive selection.
For example, mutations at Spike position 484 in SARS-CoV-2---enabling immune evasion by altering the antibody binding surface---emerged independently across multiple lineages: E484K (glutamic acid to lysine) in Beta, Gamma, and Mu, and E484A (glutamic acid to alanine) in Omicron, each descending from different ancestral branches of the phylogenetic tree. This repeated independent emergence at the same position provides strong evidence of positive selection.
Importantly, Layer~A does not impose one universal biological mechanism: the annotation can reflect immune evasion, receptor-binding adaptation, diversifying selection, or recurrent lineage-level emergence depending on the pathogen and source study.
These positions are used as \textbf{True Positives}, evaluating whether detection methods correctly identify them as hotspots.

\textbf{Identification criteria.}
For each pathogen, Layer~A positions were compiled from published studies (Table~\ref{tab:gt_positions}, rightmost column). The construction is operationally heterogeneous across the 11 pathogens: SARS-CoV-2 and H3N2 use phylogenetically confirmed convergent evolution sites ($\geq$3 independent lineages); Norovirus uses epochal-variant positions; HCV uses the hypervariable regions HVR1/HVR2 (diversifying selection rather than convergence); and the remaining seven pathogens use literature-reported immune-escape position lists. No single computational criterion is applied uniformly---the framework adopts the per-pathogen literature-curated list to maximize biological fidelity, accepting in exchange that ``Layer~A'' denotes a heterogeneous adaptive-selection union (convergent evolution $\cup$ immune escape $\cup$ HVR membership) rather than a single algorithmic construct. This heterogeneity is revisited in Chapter~\ref{ch:discussion} where its quantitative impact on the headline $\omega^2 = 0.296$ is discussed; for the purposes of this chapter the operational definition is: \emph{Layer~A is the union of literature-curated positive/diversifying-selection positions per the references in Table~\ref{tab:gt_positions}}.
For example, SARS-CoV-2 Layer~A positions (e.g., E484, N501, L452, K417) were identified by Carabelli et al.~\cite{carabelli2023convergent} through phylogenetic analysis showing independent emergence across Alpha, Beta, Gamma, Delta, and Omicron lineages.
No computational tool was applied to identify new convergent evolution positions; all positions were taken directly from the cited literature.

\subsubsection{Layer B: Constrained Purifying Selection}

Structurally and functionally essential sites under purifying selection are identified as true negatives.
These are positions where mutations are actively \textit{removed} from the viral population because they destroy protein folding, disrupt membrane fusion, or abolish receptor binding---any mutation here is lethal or severely deleterious to the virus.
Unlike Layer~A (which identifies positions under \textit{positive} selection, where change is advantageous), Layer~B identifies positions under \textit{purifying} selection (where change is harmful).
These positions exhibit low mutation rates due to loss-of-function consequences and serve as a curated constrained-negative reference set: in MCC they are subsumed under the standard true-negative pool (all non-Layer-A positions; Section~\ref{subsec:9pathogen_metrics}), while the constrained false positive rate (constrained FPR) is computed against Layer~B specifically.
When detection methods incorrectly identify Layer~B positions as hotspots, they are counted as false positives within the constrained FPR computation.

\textbf{Identification criteria.}
Layer~B positions were identified based on two types of evidence from the cited literature: (1)~structural essentiality---positions within domains whose disruption abolishes protein function (e.g., fusion peptide, HR1/HR2 heptad repeats, structural disulfide-bonding cysteine residues), as determined by published mutagenesis or structural studies; and (2)~evolutionary conservation---positions exhibiting dN/dS $< 0.5$ across diverse viral isolates, indicating strong purifying selection against amino acid changes. The 0.5 threshold follows the convention used in comparative genomics for identifying purifying selection~\cite{yang2007paml}, where $\omega < 1$ indicates purifying selection and values below 0.5 represent particularly strong functional constraint.
For SARS-CoV-2, the 154 Layer~B positions include the HR1 (AA 912--984) and HR2 (AA 1,163--1,213) heptad repeats and structural cysteine residues, identified from the cryo-EM structure~\cite{wrapp2020cryo} and conservation analysis.
As with Layer~A, no de novo computational identification was performed; all positions were curated from published sources.

\subsubsection{Layer C: DMS Experimental Validation}

For six pathogens with available Deep Mutational Scanning (DMS) experimental data---SARS-CoV-2~\cite{dadonaite2023dms}, H3N2, HIV-1, RSV, Rabies, and EV-A71---experimentally measured fitness effects are used as ground truth.
Unlike Layer~A (which infers importance from evolutionary patterns) and Layer~B (which infers importance from structural conservation), Layer~C provides direct experimental measurements: in DMS experiments, every possible single amino acid substitution at every position is physically constructed in the laboratory, and its effect on a quantifiable phenotype (e.g., cell entry efficiency for SARS-CoV-2, replicative fitness for H3N2) is measured.
Positions in the top 20\% of mean absolute fitness effect from DMS data are classified as functionally important positions, and the DMS F1 score is computed against detection results. The 20\% threshold was chosen to balance sensitivity and specificity: it captures positions with substantial fitness effects while maintaining sufficient contrast with non-functional positions. Sensitivity analysis (Section~\ref{subsec:dms_evolution_correlation}) checked Layer~C rankings at thresholds 10\%, 20\%, and 30\% (4 pathogens with available DMS at the time of the sensitivity sweep: SARS-CoV-2, H3N2, HIV-1, RSV) and found rankings stable across that range.
This layer is particularly valuable because it captures functional importance that may not yet be visible in evolutionary data (a position may have strong fitness effects but insufficient time for convergent evolution to manifest).
The DMS phenotypes measured differ by pathogen: cell entry fitness for SARS-CoV-2~\cite{dadonaite2024nature} and RSV~\cite{simonich2026rsvdms}, viral replicative fitness for H3N2~\cite{lee2018h3n2dms}, viral growth in cell culture for HIV-1~\cite{haddox2018hiv1dms}, G-protein-mediated cell entry for Rabies~\cite{aditham2025rabies}, and replicative fitness for EV-A71~\cite{bakhache2025eva71}. Table~\ref{tab:dms_preprocessing} summarises the per-pathogen DMS preprocessing conventions adopted in this benchmark; cell-line and replicate-aggregation choices follow the released-source defaults except where noted.

\begin{table}[htbp]
\centering
\caption{DMS preprocessing conventions per Layer~C source. ``Mean abs.\ effect'' refers to the per-position mean absolute fitness effect across all amino-acid substitutions, used as the score from which the top-20\% threshold defines Layer~C positives. All sources used the released processed-fitness tables; raw read filtering and replicate aggregation follow each source's published default unless noted.}
\label{tab:dms_preprocessing}
\small
\adjustbox{max width=\textwidth}{
\begin{tabular}{llllll}
\toprule
\textbf{Pathogen} & \textbf{Phenotype} & \textbf{Assay system} & \textbf{Replicates aggregated} & \textbf{Score used} & \textbf{Threshold} \\
\midrule
SARS-CoV-2 & Cell entry fitness & Pseudovirus (HEK293T) & Mean across 2 replicates & Mean abs.\ effect & Top 20\% \\
H3N2       & Replicative fitness & MDCK cell culture       & Mean across 2 replicates & Mean abs.\ effect & Top 20\% \\
HIV-1      & Viral growth        & TZM-bl cell culture     & Mean across 2 replicates & Mean abs.\ effect & Top 20\% \\
RSV        & Cell entry fitness  & Pseudovirus             & Mean across 2 replicates & Mean abs.\ effect & Top 20\% \\
Rabies     & G-protein cell entry & Pseudovirus            & Mean across 2 replicates & Mean abs.\ effect & Top 20\% \\
EV-A71     & Replicative fitness  & RD cell culture        & Mean across 2 replicates & Mean abs.\ effect & Top 20\% \\
\bottomrule
\end{tabular}
}
\begin{minipage}{\textwidth}
\smallskip
\footnotesize Sensitivity to the threshold (sweep over $10\%$, $20\%$, $30\%$) was checked in Section~\ref{subsec:dms_evolution_correlation}.
\end{minipage}
\end{table}

\textbf{Layer A/B overlap handling.}
In RSV (positions 148, 152) and Influenza B (positions 141, 194), overlapping positions exist between Layer A and Layer B.
This occurs because these residues are simultaneously under immune evasion pressure (positive selection) and structural constraints (purifying selection), representing dual-selection sites.
In MCC evaluation, these positions are prioritized as Layer A (positive) and excluded from the Layer B (negative) set.
Therefore, the effective Layer B count is 16 for RSV (18 $-$ 2) and 8 for Influenza B (10 $-$ 2) positions.

\textbf{Layer A/C conflict handling.}
Layer~A (literature-curated adaptive/diversifying evidence) and Layer~C (DMS fitness) can disagree at individual positions: a residue may be literature-supported as adaptive or immune-relevant (Layer~A positive) but have low DMS fitness cost (Layer~C negative), or vice versa. This does not require a merged classification decision because the two layers feed \textit{different} evaluation metrics under non-overlapping scopes:
(1)~\textbf{Layer~A $\to$ MCC}: Layer~A positions are treated as positives and all other positions as negatives; Layer~C membership does not affect the MCC calculation.
(2)~\textbf{Layer~C $\to$ DMS F1}: Layer~C top-20\% positions are treated as positives (for 6 pathogens with DMS); Layer~A membership does not affect the DMS F1 calculation.
(3)~\textbf{Interpretation rule}: when the two metrics rank scoring methods differently for the same pathogen, both rankings are reported without collapsing them into a single verdict; the disagreement is itself a result (Section~\ref{subsec:dms_evolution_correlation}, Table~\ref{tab:layer_a_vs_c}) reflecting that evolutionary and functional importance are non-identical dimensions of hotspot-ness (as stated in the three-orthogonal-dimensions rationale above).
Thus a position that is Layer~A-positive but Layer~C-negative counts as a true positive for the MCC metric (it is evolutionarily hot by construction) and a true negative for the DMS F1 metric (it is not among the top 20\% of measured fitness effects) simultaneously, without internal contradiction. This is consistent with the design choice that no single universal ``hotspot'' label is asserted; rather, the benchmark reports performance against each operationalization of hotspot-ness independently.

Table~\ref{tab:gt_positions} summarizes the amino acid position counts and sources for each pathogen's 3-layer ground truth.

\begin{table}[htbp]
\centering
\caption{3-layer ground truth position counts per pathogen. Counts are evaluated on the benchmark alignment coordinate system; literature residue labels are retained in the source annotations where biologically meaningful.}
\label{tab:gt_positions}
\small
\adjustbox{max width=\textwidth}{
\begin{tabular}{llrrrl}
\toprule
\textbf{Pathogen} & \textbf{Surface protein} & \textbf{Layer A} & \textbf{Layer B} & \textbf{Layer C} & \textbf{Layer A source} \\
\midrule
SARS-CoV-2 & Spike & 25 & 154 & 252 & \cite{cao2023ba286,carabelli2023convergent} \\
H3N2 & HA & 25 & 10 & 102 & \cite{koel2013substitutions} \\
Norovirus & VP1 & 15 & 229 & --- & \cite{lindesmith2012,lindesmith2013} \\
HIV-1 & Env gp120 & 23 & 23 & 100 & \cite{moore2015williamson} \\
Dengue & E protein & 24 & 24 & --- & \cite{twiddy2002dengue,holmes2003dengue} \\
RSV & F protein & 20 & 18 & 101 & \cite{mas2018rsv,sullender2000rsv} \\
Influenza B & HA & 17 & 10 & --- & \cite{ni2013influenzab,wang2008influenzab} \\
MERS & Spike & 9 & 110 & --- & \cite{kim2016mers,tang2014mers} \\
HCV & E2 & 54 & 0$^\ddagger$ & --- & HVR1/HVR2 hypervariable regions \\
Rabies & G protein & 39 & 42 & 87 & \cite{benmansour1991rabies,aditham2025rabies} \\
EV-A71 & VP1 & 27 & 29 & 55 & \cite{huang2015eva71,bakhache2025eva71} \\
\bottomrule
\end{tabular}
}
\begin{minipage}{\textwidth}
\smallskip
\footnotesize Layer C (DMS) sources: SARS-CoV-2~\cite{dadonaite2023dms}, H3N2~\cite{lee2018h3n2dms}, HIV-1~\cite{haddox2018hiv1dms}, RSV~\cite{simonich2026rsvdms}, Rabies~\cite{aditham2025rabies}, EV-A71~\cite{bakhache2025eva71}.
``---'' indicates that no DMS data is available for that pathogen.
$^\ddagger$HCV E2 literature coordinates are commonly reported in polyprotein numbering, whereas the benchmark feature matrix is indexed on the 340-AA E2 alignment. HVR1/HVR2 were mapped into the evaluated E2 alignment, but CD81-binding regions reported as AA 418--535 in the literature were not converted into a high-confidence alignment-internal constrained set in this release; HCV therefore has effective Layer~B = 0. HCV-specific conclusions should be read with this coordinate-handling caveat, and Chapter~\ref{ch:results} reports HCV-exclusion sensitivity for the headline interaction.

\textbf{HIV-1 Layer~A count reconciliation}: this table reports the literature-curated Layer~A count of \textbf{23} positions (Moore \& Williamson 2015 baseline) which is the canonical narrative reference. The downstream operational alignment-QC count of \textbf{26} positions (used by AUC computation in Table~\ref{tab:gt_heterogeneity} and the positive-rate footnote in Table~\ref{tab:design_characteristics}) adds three positions identified during alignment quality control after the literature scan; the difference does not affect any headline conclusion (the scoring~$\times$~pathogen interaction is computed identically under both counts because Layer~A appears as a binary label, not a count, in the per-position regression).

Note that Layer~A selection criteria vary across pathogens: SARS-CoV-2 and H3N2 use phylogenetically confirmed convergent evolution sites, Norovirus uses epochal variant positions, HCV uses hypervariable regions (HVR1/HVR2), and the remaining pathogens use literature-reported immune evasion positions. This heterogeneity is a limitation of the current ground truth construction (discussed in Chapter~\ref{ch:discussion}).

\textbf{Inclusion of pathogens with incomplete ground truth layers.}
Six pathogens (SARS-CoV-2, H3N2, HIV-1, RSV, Rabies, EV-A71) have all three layers available.
Norovirus, Dengue, Influenza~B, MERS, and HCV lack Layer~C (DMS data).
These pathogens are nevertheless included because (1)~the primary evaluation metric (MCC) is computed using only Layer~A (positive selection sites as positives, all other positions as negatives), so the core evaluation is valid without Layers~B/C;
(2)~restricting the benchmark to six pathogens would provide insufficient statistical power for the cross-pathogen analyses (ANOVA, Friedman, LOPO) in Chapter~\ref{ch:results};
and (3)~including pathogens with diverse evolutionary dynamics and varying ground truth availability more accurately reflects real-world conditions where comprehensive ground truth is rarely available for all pathogens.
\end{minipage}
\end{table}

Figure~\ref{fig:gt_distribution} visualizes the spatial distribution of Layer~A and Layer~B positions across all 11 pathogens, illustrating the diversity in ground truth density, protein length, and the balance between positive and negative annotations.

\begin{figure}[htbp]
\centering
\includegraphics[width=\textwidth]{gt_distribution_11pathogens.png}
\caption{Spatial distribution of 3-layer ground truth positions across 11 pathogens. Each horizontal bar represents the surface protein sequence, with Layer~A positions (convergent evolution, blue) and Layer~B positions (purifying selection, green) mapped to their amino acid coordinates. Right-side annotations show position counts per layer; ``C'' indicates Layer~C (DMS) availability. Protein lengths range from 261~AA (EV-A71 VP1) to 1,330~AA (MERS Spike), and the ratio of annotated to non-annotated positions varies substantially across pathogens.}
\label{fig:gt_distribution}
\end{figure}

\subsection{Stage~1: SARS-CoV-2 In-Depth Analysis}
\label{subsec:stage1}

This subsection describes the Stage~1 experimental design, which focuses on SARS-CoV-2 as the anchor pathogen.
Stage~1 serves three purposes: (1)~validating the framework components on the pathogen with the richest ground truth data, (2)~comparing MutClust variants to establish the detection baseline, and (3)~analyzing temporal scoring behavior.
The ground truth, detection methods, evaluation metric, and scoring variants used in Stage~1 are described below.

\subsubsection{Ground Truth Generation}
\label{subsec:stage1_gt}

Ground truth is generated from three sources.

\textbf{DMS-based functional sites.}\phantomsection\label{subsec:dms_ground_truth}
The DMS fitness data are derived from the full-length Spike protein characterization experiment by Dadonaite et al.~\cite{dadonaite2023dms, dadonaite2024nature}.
Pseudovirus cell entry fitness measurements are used, as this is the most directly relevant phenotype capturing the overall functional impact of mutations on viral entry.
Amino acid positions in the top 20\% of mean absolute fitness effect are selected as functional sites and mapped to nucleotide coordinates through codon-aware coordinate transformation:
\begin{equation}
\text{nuc\_positions} = \{g_{\text{start}} + (p_{\text{aa}} - 1) \times 3 + k : k \in \{0, 1, 2\}\}
\label{eq:coordinate_transform}
\end{equation}
where $g_{\text{start}} = 21{,}563$ is the start position of the Spike gene in the Wuhan-Hu-1 reference genome.
This provides an experimental ground truth independent of the detection methods under evaluation, thereby avoiding the circularity of using bootstrap p-values for validation.

\textbf{Synthetic data (controlled signal recovery test).}
The synthetic benchmark was designed as a controlled signal recovery test that evaluates how accurately each detection method recovers artificially inserted signals at known positions.
Synthetic H-scores are generated mimicking SARS-CoV-2 patterns, with hotspot regions inserted at known functional sites.
High H-scores (0.8--1.2) are assigned to Core Candidate Mutation (CCM) seed positions (AA 484, 501, 417, 452---the four Spike positions most frequently mutated in major SARS-CoV-2 variants), moderate scores (0.05--0.3) to surrounding functional regions, and low noise (0.0--0.02) to background positions.
Scores within each category are drawn from uniform distributions: $\mathcal{U}(0.8, 1.2)$ for seed positions, $\mathcal{U}(0.05, 0.3)$ for surrounding functional regions (defined as positions within 50 positions of each seed), and $\mathcal{U}(0.0, 0.02)$ for background. A single synthetic genome of 29,903 positions is generated per trial.
Since the ground truth of synthetic data is known by design, this test evaluates each method's signal recovery capability but does not directly measure actual biological hotspot detection performance.
Therefore, DMS-based validation (Section~\ref{subsec:dms_ground_truth}) and convergent evolution evaluation (Section~\ref{sec:multi_gt}) provide independent biological validation, confirming whether the conclusions from the synthetic benchmark hold on real data.

\textbf{Annotated functional regions.}
The annotated functional regions used as ground truth are presented in Table~\ref{tab:ground_truth}, with region boundaries based on the cryo-EM structure of the prefusion Spike trimer~\cite{wrapp2020cryo} and the N-terminal domain (NTD) antigenic supersite mapping by McCallum et al.~\cite{mccallum2021ntd}.

\begin{table}[htbp]
\centering
\caption{Functional regions of the Spike protein (ground truth)}
\label{tab:ground_truth}
\adjustbox{max width=\textwidth}{
\begin{tabular}{llr}
\toprule
\textbf{Functional Region} & \textbf{Amino Acid (AA) Positions} & \textbf{Nucleotide Count} \\
\midrule
Receptor Binding Domain (RBD) & AA 319--541 & 669 \\
NTD Antigenic Supersite & AA 14--20, 140--158, 245--264 & 138 \\
Furin Cleavage Site (S1/S2) & AA 681--685 & 15 \\
S2' Cleavage Site & AA 815--816 & 6 \\
Fusion Peptide & AA 816--833 & 54$^*$ \\
Heptad Repeat 1 (HR1) & AA 912--984 & 219 \\
Heptad Repeat 2 (HR2) & AA 1,163--1,213 & 153 \\
\midrule
\textbf{Total (deduplicated)} & & \textbf{1,251}$^\dagger$ \\
\bottomrule
\end{tabular}
}
\begin{minipage}{\textwidth}
\smallskip
\footnotesize $^*$Overlaps with the S2' cleavage site at AA 816 (3 nucleotides shared).
$^\dagger$The sum of individual regions is 1,254 nt; after removing the 3 nucleotide (nt) overlap at the S2'/fusion peptide boundary (AA 816), 1,251 unique nucleotide positions constitute the ground truth.
\end{minipage}
\end{table}

The genome-wide distribution of these functional regions is visualized in Figure~\ref{fig:mutbench_hotspot_map}, which maps the seven annotated regions to their nucleotide coordinates on the full SARS-CoV-2 genome (29,903 positions). The ground truth positions are concentrated within the Spike gene (nucleotides 21,563--25,384), while the remainder of the genome shows minimal functional annotation, confirming the appropriateness of Spike-focused analysis.

\begin{figure}[htbp]
\centering
\includegraphics[width=\textwidth]{genome_hotspot_map.png}
\caption{Genome-wide hotspot map of the SARS-CoV-2 Spike protein. The seven annotated functional regions used as ground truth (RBD, NTD antigenic supersite, furin cleavage site, S2' cleavage site, fusion peptide, HR1, HR2) are mapped to nucleotide coordinates, encompassing 1,251 unique Spike nucleotide positions.}
\label{fig:mutbench_hotspot_map}
\end{figure}

\subsubsection{MutClust Variant Comparison}
\label{subsec:variants}

Among the 14 detection families used in the full benchmark, MutClust~\cite{youn2025mutclust} is the only virus-specific hotspot detection tool and serves as the starting point for this study.
MutClust scores each genomic position using the H-score:
\begin{equation}
H_i = \log_2(P_i \times E_i \times 100 + 1)
\label{eq:mutbench_hscore}
\end{equation}
where $P_i$ is the sum of mutation frequencies at position $i$ and $E_i$ is the Shannon entropy of mutation types:\phantomsection\label{subsec:mutbench_hscore}
\begin{equation}
E_i = -\sum_{j=1}^{k} p_{ij} \log_2 p_{ij}
\label{eq:shannon_entropy}
\end{equation}
Hotspot regions are then identified via adaptive DBSCAN clustering with $\varepsilon = \gamma \times H_i$, where the scaling constant $\gamma = 0.5$ is fixed across all pathogens (sensitivity to $\gamma \in \{0.25, 0.5, 1.0\}$ shifts the per-pathogen MCC by $\pm 0.01$ on average, well within the subsampling band reported in Section~\ref{sec:robustness_validation}).\phantomsection\label{eq:expand_cluster}
The internal DBSCAN call uses \texttt{min\_samples} = 3, and the HDBSCAN baseline (v3) and HDBSCAN-Auto detector use \texttt{min\_cluster\_size} = 5 with \texttt{min\_samples} = 1 (default \texttt{cluster\_selection\_method = 'eom'}); these values are pinned in the released configuration files described in the \textbf{Code and Data Availability} subsection at the end of this chapter.

Five variants were tested in Stage~1: the original implementation (v1, with a cumulative decay bug), the bug-fixed version (v2a), a dN/dS-filtered version (v2b), a CCM-seed + HDBSCAN hybrid (v2c), and a pure HDBSCAN baseline (v3).
MutClust-Hybrid (v2c) achieved the highest hotspot-score (0.778) in Stage~1 validation.
In the Stage~2 cross-pathogen benchmark, only MutClust-Orig (v1) is retained as one of the 14 detection families, for two reasons.
First, Stage~2 shifts focus from algorithm comparison to information-type analysis, requiring each detection family to accept \textit{any} of the 20 scoring types as input.
MutClust-Hybrid (v2c) achieves its advantage through H-score-specific domain knowledge (CCM seed selection assumes H-score-like distributions), making it incompatible with non-H-score inputs such as ESM-2 LLR or pLDDT---its CCM seed logic becomes meaningless for these scoring types.
MutClust-Orig (v1), in contrast, applies adaptive DBSCAN to any input score array without scoring-specific assumptions, ensuring fair comparison across all 20 scoring types.
Second, using the published version~\cite{youn2025mutclust} as the baseline ensures reproducibility and represents the prior state-of-the-art as reported in the literature.\phantomsection\label{subsec:improvements}

\subsubsection{Hotspot-Score Composite Metric}
\label{subsec:hotspot_score}

The hotspot-score is a composite evaluation metric designed for Stage~1 that integrates three complementary perspectives:
\begin{equation}
\text{hotspot-score} = \frac{\text{recall} + \text{precision} + \text{stability}}{3}
\label{eq:hotspot_score}
\end{equation}
where stability is the mean Jaccard similarity ($J = |A \cap B| / |A \cup B|$) across 100 random 80\% subsampling iterations.
In each iteration, a Bernoulli mask independently zeroes out approximately 20\% of genomic positions, the detection method is rerun on the perturbed scores, and the Jaccard similarity with the full-data result is computed.
This metric balances detection accuracy (recall, precision) with reproducibility (stability), reflecting the practical requirement that methods in viral surveillance should produce consistent results under data perturbation.
Stage~2 uses MCC instead of hotspot-score (see Section~\ref{subsec:9pathogen_metrics} for the rationale).

\subsubsection{Temporal Scoring Variants}
\label{subsec:alternative_scoring}

The following scoring variants were developed specifically for the Stage~1 SARS-CoV-2 temporal analysis and are distinct from the 20 Stage~2 scoring types described in Section~\ref{subsec:9pathogen_scoring}.
As H-score saturation was identified in the temporal validation (Section~\ref{sec:temporal}), alternative scoring methods were designed for complementary evaluation against various ground truth definitions.
These methods assign a score to each genomic position; top positions are selected via thresholding and compared against multiple ground truths.
Table~\ref{tab:alternative_scoring} summarizes these Stage~1 variants.

\begin{table}[htbp]
\centering
\caption{Stage~1 temporal scoring variants used only for SARS-CoV-2 analysis}
\label{tab:alternative_scoring}
\small
\adjustbox{max width=\textwidth}{
\begin{tabular}{lll}
\toprule
\textbf{Name} & \textbf{Formula} & \textbf{Description} \\
\midrule
Original H-score & Eq.~\ref{eq:mutbench_hscore} & MutClust's original H-score \\
Saturated H-score & (same, on 2024--25 data) & H-score on contemporary consensus alignment \\
CH-score & $\log_2(P_i^{\text{cons}} \times E_i \times 100 + 1)$ & Consensus-based; removes reference dependency \\
E-score & $-\sum_{a} f_{a,i} \log_2 f_{a,i}$ & Shannon entropy alone \\
TC-freq & $|\text{freq}_{t_2} - \text{freq}_{t_1}| / (\text{freq}_{t_1} + \epsilon)$ & Temporal frequency change rate \\
TC-freq-v2 & $\text{TC-freq}_i \times \text{freq}_{t_2}(i)$ & TC-freq with absolute frequency weighting \\
TC-freq-v3 & $\max_q \text{TC-freq}_{q}$ & Maximum quarterly frequency change \\
dN/dS-weighted CH & $\text{CH}_i \times (dN/dS)_i$ & CH-score weighted by position-level dN/dS \\
\bottomrule
\end{tabular}
}
\end{table}

\noindent For each scoring method, thresholds at the top 5\%, 10\%, 20\%, and median are applied to generate sets of ``detected'' positions, and F1, enrichment ratio, and Cohen's $d$ effect size are computed for each ground truth.

%% ============================================================
\subsection{Stage~2: 11-Pathogen Cross-Pathogen Benchmark}
\label{sec:stage2}

This subsection describes the Stage~2 experimental design, which extends the evaluation from SARS-CoV-2 to all 11 pathogens.
Stage~2 systematically evaluates 20 scoring types $\times$ 39 detection parameter variants (from 14 families) $\times$ 11 pathogens = 8,580 combinations to identify the dominant factors influencing hotspot detection performance.

\subsubsection{Scoring Types}
\label{subsec:9pathogen_scoring}
\label{sec:scoring_detection}

Twenty scoring types organized into six categories are defined to assign per-position scores from diverse information sources (Table~\ref{tab:scoring_types}).
Each category captures a distinct biological or computational perspective on mutational importance, ensuring that the benchmark evaluates detection methods across a broad range of input signals.

\begin{table}[htbp]
\centering
\caption{Twenty scoring types used in the MutBench Stage~2 benchmark (6 categories)}
\label{tab:scoring_types}
\small
\adjustbox{max width=\textwidth}{
\begin{tabular}{llll}
\toprule
\textbf{Category} & \textbf{Scoring type} & \textbf{Information source} & \textbf{Description} \\
\midrule
\multirow{3}{*}{Frequency-based (3)} & freq & MSA mutation frequency & Position-level mutation frequency \\
 & rare\_freq & MSA rare variant frequency & Frequency of minor alleles below 5\% threshold (low-frequency variants signalling emerging adaptation) \\
 & freq\_with\_indel & MSA frequency incl.\ indels & Mutation frequency including insertion/deletion events \\
\midrule
\multirow{3}{*}{MSA-derived (3)} & entropy & MSA Shannon entropy & Per-position diversity ($-\sum p_j \log_2 p_j$) \\
 & entropy\_with\_indel & MSA entropy incl.\ indels & Shannon entropy including gap characters \\
 & homoplasy & Phylogenetic tree & Convergent evolution count on reconstructed phylogeny \\
\midrule
\multirow{4}{*}{Phylogenetic (4)} & dN/dS proxy & Codon alignment & Nei--Gojobori nonsynonymous/synonymous ratio proxy \\
 & FUBAR prob & Codon model (HyPhy) & Posterior probability of positive selection \\
 & FUBAR pos\_sel & Codon model (HyPhy) & Positive selection rate ($\beta - \alpha$) \\
 & FUBAR Bayes factor & Codon model (HyPhy) & Bayes factor evidence for positive selection \\
\midrule
\multirow{3}{*}{Structural (3)} & pLDDT inverse & AlphaFold/ESMFold structure & $1 - \text{pLDDT}$; low confidence $\approx$ flexible/disordered \\
 & SASA & 3D structure & Solvent-accessible surface area (exposed residues) \\
 & Grantham distance & Physicochemical properties & Mean Grantham distance of observed substitutions \\
\midrule
\multirow{5}{*}{AI-based (5)} & ESM-2 LLR & Protein language model & Masked-marginal mutation-surprise score (LLR shorthand) \\
 & Tranception\textsuperscript{$\ast$} & Autoregressive PLM proxy & Per-position pseudo-perplexity (proxy; see implementation note below) \\
 & semantic change & ESM-2 embeddings & Cosine distance in embedding space (WT vs.\ mutant) \\
 & stability$^{\dagger}$ & Substitution-matrix proxy & BLOSUM62-based stability proxy (radical substitutions $\to$ destabilizing) \\
 & stability\_structural$^{\dagger}$ & Structure-weighted proxy & BLOSUM62 score weighted by ESMFold pLDDT (well-folded $\to$ higher impact) \\
\midrule
\multirow{2}{*}{Composite (2)} & EVEscape composite$^{\dagger}$ & EVEscape-style proxy & Multiplicative composite of ESM-2 LLR $\times$ SASA $\times$ Grantham (proxy for fitness $\times$ accessibility $\times$ dissimilarity) \\
 & EqualWeight & All-feature average & Equal-weighted average of all normalized scoring types \\
\bottomrule
\end{tabular}
}
\end{table}

\textbf{Frequency-based scores} capture the raw mutation signal from the MSA: how often each position deviates from the consensus, including rare variants and indel-inclusive counts.
These scores reflect the most direct observable signal but are susceptible to sampling bias and alignment artifacts.

\textbf{MSA-derived scores} quantify diversity patterns beyond simple frequency: Shannon entropy measures allelic diversity, while homoplasy counts independent convergent mutations on the phylogeny, providing a signal distinct from overall frequency.
Phylogenetic trees were reconstructed using IQ-TREE~2~\cite{minh2020iqtree} with automatic model selection (ModelFinder) and 1{,}000 ultrafast bootstrap replicates. Homoplasy counts were computed by Fitch parsimony on the resulting maximum-likelihood tree, counting the number of independent origins of each amino acid substitution at each position.

\textbf{Phylogenetic scores} leverage codon-level evolutionary models to estimate selection pressures.
The dN/dS proxy uses the Nei--Gojobori method, while the three FUBAR-based scores~\cite{murrell2013fubar} provide Bayesian estimates of site-specific selection from the HyPhy framework, each capturing a different aspect of the posterior distribution.
FUBAR analysis was performed using HyPhy v2.5 with the MG94$\times$REV codon model and a default posterior probability threshold of 0.9.

\textbf{Structural scores} incorporate three-dimensional protein information: pLDDT inverse from AlphaFold DB or ESMFold-derived structures identifies flexible or disordered regions, SASA measures surface exposure (relevant for immune accessibility), and Grantham distance~\cite{grantham1974amino} quantifies the physicochemical severity of observed substitutions. The per-pathogen structure source used for pLDDT and SASA computation is summarized in Table~\ref{tab:structure_sources}: SARS-CoV-2 uses the AlphaFold~DB entry, while all other pathogens use the ESMFold API.
SASA values were computed from the available predicted structures using a solvent-accessibility calculation; the main result tables report the ESMFold/Shrake--Rupley-derived release, while earlier Lee--Richards/FreeSASA pilot outputs are treated as archival.

\begin{table}[htbp]
\centering
\caption{Per-pathogen structure source for pLDDT and SASA computation. SARS-CoV-2 uses the AlphaFold~DB entry; all other pathogens use the ESMFold API for both pLDDT and SASA features. The 3D spatial-contact analysis (Chapter~\ref{ch:results}, Section~\ref{sec:structure_analysis}) uses the cryo-EM structure PDB~6VSB~\cite{wrapp2020cryo} and is restricted to SARS-CoV-2.}
\label{tab:structure_sources}
\footnotesize
\begin{tabular}{llll}
\toprule
\textbf{Pathogen} & \textbf{Target protein} & \textbf{pLDDT/SASA source} & \textbf{3D spatial-contact} \\
\midrule
SARS-CoV-2 & Spike (S) & AlphaFold~DB (AF-0000000365840314) & PDB~6VSB cryo-EM \\
H3N2 (Influenza A) & HA & ESMFold API & not analyzed \\
Influenza~B & HA & ESMFold API & not analyzed \\
HCV & E2 & ESMFold API & not analyzed \\
HIV-1 & gp120 & ESMFold API & not analyzed \\
RSV & F & ESMFold API & not analyzed \\
Norovirus & VP1 & ESMFold API & not analyzed \\
Dengue & E & ESMFold API & not analyzed \\
MERS & Spike (S) & ESMFold API & not analyzed \\
Rabies & G & ESMFold API & not analyzed \\
EV-A71 & VP1 & ESMFold API & not analyzed \\
Zika$^{\dagger}$ & E & ESMFold API & not analyzed \\
\bottomrule
\end{tabular}
\begin{minipage}{\textwidth}
\smallskip
\footnotesize $^{\dagger}$Zika is included only in the Stage~3 12-pathogen feature-ablation panel (Section~\ref{subsec:stage3_methods}); the main 11-pathogen Stage~2 benchmark excludes Zika.
\end{minipage}
\end{table}

\textbf{AI-based scores} leverage protein language models (PLMs) trained on large-scale protein sequence databases.
The ESM-2 channel is labelled ``ESM-2 LLR'' for consistency with protein-language-model variant-effect literature, but operationally it is a masked-marginal pseudo-log-likelihood / mutation-surprise score rather than a full mutation-vs-wildtype likelihood-ratio model trained for each pathogen. Tranception-style autoregressive scores predict mutational effects from evolutionary context, semantic change captures embedding-space shifts, and stability predictors estimate thermodynamic consequences of mutations.
\textit{Implementation note:} the ``Tranception'' scoring channel in this benchmark is computed using ESM-2 masked-marginal pseudo-perplexity as a lightweight proxy for the Tranception autoregressive objective~\cite{notin2022tranception}; the original Tranception model was not run directly. Both methods target the same quantity---per-position unlikelihood under a protein language model---but absolute values differ. To audit channel separability, Table~\ref{tab:tranception_esm_correlation} reports the per-pathogen Pearson and Spearman correlations between the Tranception channel and the standalone ESM-2 LLR channel on the position-level scoring vectors used in Stage~2. The two channels are highly correlated for 11 of 12 pathogens (Pearson $\rho = 0.64$--$0.84$, Spearman $\rho = 0.71$--$0.98$; all $p < 10^{-50}$ except Zika), consistent with their shared ESM-2 pseudo-perplexity basis. HIV-1 is the lone exception (Pearson $\rho = 0.04$, $p = 0.36$ n.s.; Spearman $\rho = 0.16$, $p = 7.8 \times 10^{-4}$). The HIV-1 dissociation is consistent with the high antigenic diversity and low conservation of gp120, where the relative ordering of positions diverges from the magnitude-based score even within the same pseudo-perplexity backbone. Readers comparing pathogen-level best-scoring-type results (e.g., SARS-CoV-2 best $=$ Tranception vs.\ RSV best $=$ ESM-2 LLR) should therefore treat these two channels as near-equivalent for 11 of 12 pathogens and as channel-distinguishable only for HIV-1.

\begin{table}[htbp]
\centering
\caption{Per-pathogen correlation between the Tranception channel (ESM-2 masked-marginal pseudo-perplexity proxy) and the standalone ESM-2 LLR channel. Both channels share the ESM-2 \texttt{esm2\_t33\_650M\_UR50D} pseudo-perplexity backbone but use different normalization; this table audits whether they remain operationally distinct within the benchmark. Correlations were computed on the position-level scoring vectors used in Stage~2; $p$-values are uncorrected.}
\label{tab:tranception_esm_correlation}
\footnotesize
\begin{tabular}{lrrrrr}
\toprule
\textbf{Pathogen} & \textbf{N} & \textbf{Pearson $\rho$} & \textbf{Pearson $p$} & \textbf{Spearman $\rho$} & \textbf{Spearman $p$} \\
\midrule
SARS-CoV-2     & 1259 & 0.753 & $< 10^{-200}$ & 0.963 & $< 10^{-300}$ \\
H3N2           &  543 & 0.803 & $< 10^{-120}$ & 0.965 & $< 10^{-300}$ \\
Influenza~B    &  560 & 0.825 & $< 10^{-130}$ & 0.939 & $< 10^{-260}$ \\
HCV            &  340 & 0.737 & $< 10^{-58}$  & 0.980 & $< 10^{-200}$ \\
HIV-1          &  450 & \textbf{0.043} & \textbf{0.36 (n.s.)} & \textbf{0.158} & $7.8 \times 10^{-4}$ \\
RSV            &  550 & 0.805 & $< 10^{-120}$ & 0.967 & $< 10^{-300}$ \\
Norovirus      &  536 & 0.812 & $< 10^{-120}$ & 0.973 & $< 10^{-300}$ \\
Dengue         &  490 & 0.828 & $< 10^{-120}$ & 0.946 & $< 10^{-200}$ \\
MERS           & 1330 & 0.836 & $< 10^{-300}$ & 0.963 & $< 10^{-300}$ \\
Rabies         &  490 & 0.679 & $< 10^{-65}$  & 0.949 & $< 10^{-200}$ \\
EV-A71         &  260 & 0.801 & $< 10^{-58}$  & 0.919 & $< 10^{-100}$ \\
Zika$^{\dagger}$ &  451 & 0.643 & $< 10^{-53}$  & 0.713 & $< 10^{-70}$  \\
\bottomrule
\end{tabular}
\begin{minipage}{\textwidth}
\smallskip
\footnotesize $^{\dagger}$Zika is included only in the Stage~3 12-pathogen feature-ablation panel. Underlying CSV: \texttt{results/mutbench/feature\_analysis/tranception\_esm\_correlation.csv}.
\end{minipage}
\end{table}
The \texttt{stability} channel uses a BLOSUM62-based substitution proxy: at each MSA position, $\text{stability\_score} = -\overline{\text{BLOSUM62}(\text{wt}, \text{mut})}$ weighted by mutation frequency, so that radical substitutions (low BLOSUM62) yield higher destabilization scores. The \texttt{stability\_structural} variant multiplies this score by per-residue ESMFold pLDDT/100, reflecting that mutations in well-folded regions are more likely to disrupt structure. This proxy is used in place of full physics-based ($\Delta\Delta G$) predictors such as FoldX or RoseTTAFold, which are computationally infeasible at the 11-pathogen $\times$ 39-detector scale of MutBench.
ESM-2 scores were computed using the \texttt{esm2\_t33\_650M\_UR50D} checkpoint (33 layers, 650M parameters).

\textbf{Composite scores} combine multiple information sources: EVEscape integrates evolutionary fitness, structural accessibility, and physicochemical dissimilarity, while EqualWeight provides a naive baseline averaging all normalized scores.
The \texttt{EVEscape\_composite} channel implements an EVEscape-style multiplicative combination~\cite{thadani2023evescape}: $\text{composite} = \text{ESM-2}_\text{norm} \times \text{SASA}_\text{norm} \times \text{Grantham}_\text{norm}$, where each component is min-max normalized to $[0, 1]$. This substitutes the ESM-2 masked-marginal score for the original EVE fitness term, since EVE training is not feasible across all 11 pathogens within the present scope; the structural accessibility (SASA) and physicochemical dissimilarity (Grantham) terms follow the original EVEscape construction.
Each scoring type was $z$-score normalized (zero mean, unit variance) across positions before averaging.
A directional sign rule was applied so that, prior to averaging, each feature is multiplied by $\text{sign}(\hat{s}_f)$, where $\hat{s}_f \in \{-1,+1\}$ is the empirically observed sign of the feature's per-position correlation with Layer~A on the training pathogens (e.g., \texttt{pLDDT inverse} and \texttt{rare\_freq} enter with $+1$; conserved-region indicators like raw pLDDT would enter with $-1$). This keeps EqualWeight a deterministic, label-light fusion: all magnitudes are equal-weighted, but every feature is oriented toward ``higher score $=$ more hotspot-like'' before $z$-score averaging, preventing destructive cancellation between features whose raw scales point in opposite directions.

\subsubsection{Detection Method Families}
\label{subsec:9pathogen_detection}

A total of 39 detection methods belonging to 14 algorithm families are evaluated.
Each detection method identifies hotspot regions from the per-position score profile assigned by a scoring type.
Two additional ensemble variants (majority vote of 4 threshold-based sub-detectors) were initially considered but excluded from the benchmark because their outputs are deterministic combinations of other methods already in the benchmark, violating the independence assumption required for fair comparison.

\begin{table}[htbp]
\centering
\caption{39 detection methods evaluated in the 11-pathogen benchmark (14 families)}
\label{tab:detection_families}
\small
\adjustbox{max width=\textwidth}{
\begin{tabular}{llr}
\toprule
\textbf{Family} & \textbf{Method description} & \textbf{Variants} \\
\midrule
Wavelet & Multi-resolution peak detection via wavelet transform & 3 \\
KDE & Hotspot detection via kernel density estimation & 3 \\
SlidingTest & Sliding window statistical test & 2 \\
GradPeak & Gradient peak detection on score profile & 2 \\
LocalContrast & Local contrast (local--global ratio) based detection & 3 \\
CUSUM & Cumulative sum change-point detection & 3 \\
AdaptiveKnee & Adaptive curvature-based threshold & 1 \\
ScoreDBSCAN & Score-based density-space clustering & 2 \\
Spectral & Spectral analysis-based periodic pattern detection & 2 \\
Bayes & Bayesian change-point detection & 2 \\
\midrule
AdaptWin & Adaptive window size detection & 2 \\
FreqThresh & Frequency percentile threshold (top-$k$\%) & 4 \\
SWAN & Sliding window adaptive normalization & 9 \\
MutClust-Orig & MutClust original algorithm & 1 \\
\midrule
\textbf{Total} & & \textbf{39} \\
\bottomrule
\end{tabular}
}
\end{table}

The 14 families are organized into 5 higher-level categories reflecting their algorithmic approach to converting continuous score profiles into discrete hotspot calls.
Supervised machine learning methods (e.g., Random Forest or gradient boosting classifiers) were not included because they require labeled training data at the position level, and with only 11 pathogens---each having 9--54 Layer~A positives---the available ground truth is insufficient for training generalizable classifiers without severe overfitting.
Since all methods operate on the same one-dimensional score array, the detection step serves as a necessary score-to-call conversion; the ANOVA results confirm that the detection family contributes only $\omega^2 = 0.013$ to performance variance, indicating that the choice of conversion method is a minor factor compared to the choice of input information.

The 39 methods arise from 14 algorithm families, each instantiated with different parameter settings (Table~\ref{tab:detection_variants}).
For example, the Wavelet family is tested with three peak detection thresholds ($t = 1.0, 1.5, 2.0$ standard deviations above the CWT coefficient mean), and FreqThresh is tested with four percentile cutoffs (top 1\%, 3\%, 5\%, 10\%).
This multi-parameter design ensures that the benchmark evaluates each algorithm family across a range of sensitivity--specificity trade-offs rather than at a single, potentially suboptimal, operating point.

\begin{table}[htbp]
\centering
\caption{Parameter variants for each detection method family (39 methods total)}
\label{tab:detection_variants}
\small
\adjustbox{max width=\textwidth}{
\begin{tabular}{llrl}
\toprule
\textbf{Family} & \textbf{Parameter} & \textbf{Variants} & \textbf{Values} \\
\midrule
Wavelet & CWT peak threshold $t$ (std.\ dev.) & 3 & $t = 1.0, 1.5, 2.0$ \\
KDE & Density percentile $p$ & 3 & $p = 75, 80, 85$ \\
SlidingTest & Test-level $p$ (per-window) & 2 & $p = 0.01, 0.05$ \\
GradPeak & Gradient percentile $p$ & 2 & $p = 0.2, 0.5$ \\
LocalContrast & Local $z$-score threshold & 3 & $z = 1.0, 1.5, 2.0$ \\
CUSUM & Drift $d$ $\times$ threshold $h$ & 3 & $(d{=}0.3, h{=}3.0), (d{=}0.5, h{=}4.0), (d{=}0.5, h{=}5.0)$ \\
AdaptiveKnee & (parameter-free) & 1 & Automatic knee detection \\
ScoreDBSCAN & Epsilon \texttt{ep} & 2 & $\texttt{ep} = 8, 15$ (\texttt{min\_samples} = 3 throughout) \\
Spectral & Peak $z$-score threshold & 2 & $z = 1.5, 2.0$ \\
Bayes & Prior probability & 2 & prior $= 0.10, 0.15$; Gaussian likelihood $\mathcal{N}(\mu_c, \sigma_c^2)$ per class fit by maximum likelihood, posterior $> 0.5$ \\
AdaptWin & Seed percentile $s$ & 2 & $s = 85, 90$ \\
FreqThresh & Top percentile $p$ & 4 & $p = 80, 85, 90, 95$ (equivalently top $20, 15, 10, 5\%$) \\
SWAN & Window $w$ $\times$ $k$-score $k$ & 9 & $w \in \{10, 20, 50\} \times k \in \{1.0, 1.5, 2.0\}$ \\
MutClust-Orig & Gap parameter $g$ & 1 & $g = 10$ (CCM~+ adaptive DBSCAN) \\
\midrule
\textbf{Total} & & \textbf{39} & \\
\bottomrule
\end{tabular}
}
\begin{minipage}{\textwidth}
\smallskip
\footnotesize Parameter names and values match the detector strings recorded in \texttt{stage3\_full\_results.csv} (column \texttt{detector}) and are authoritative for reproduction; older design-phase drafts used different parameter names.
\end{minipage}
\end{table}

The number of parameter variants per family ranges from 1 (AdaptiveKnee, MutClust-Orig) to 9 (SWAN), reflecting the inherent parametric complexity of each algorithm. Families with more variants have a higher probability of containing a favorable parameter setting by chance; however, the ANOVA uses detection \textit{family} (14 levels) rather than individual variants (39 levels) as a factor, partially mitigating this bias.

\noindent The 14 families are grouped into five categories below; each family is summarized in Table~\ref{tab:detection_families}.

\textbf{Threshold-based} (FreqThresh, AdaptiveKnee, AdaptWin).
These methods select hotspots by applying a score cutoff: FreqThresh uses a fixed percentile threshold (top $k$\%), AdaptiveKnee automatically determines the threshold via knee-point detection on the sorted score curve, and AdaptWin extends seed positions above a percentile threshold bidirectionally to adjacent high-scoring positions.
Frequency-based thresholding is widely used in viral genomic surveillance; for example, Mercatelli and Giorgi~\cite{mercatelli2020geographic} applied mutation frequency thresholds to identify recurrent SARS-CoV-2 mutations across geographic regions.

\textbf{Signal processing} (Wavelet, Spectral, CUSUM).
Transform-based methods that convert the score profile into an alternative representation to identify anomalous regions: Wavelet applies continuous wavelet transform with a Mexican hat kernel across multiple scales, Spectral uses FFT-based high-pass filtering to isolate local deviations, and CUSUM applies cumulative sum control charts to detect change points in the $z$-normalized profile.
Wavelet-based approaches have been applied to genomic sequence analysis, including Anastassiou~\cite{anastassiou2001genomic} who demonstrated wavelet transforms for detecting periodicities in DNA sequences, and Arneodo et al.~\cite{arneodo2011wavelet} who used multi-scale wavelet analysis to characterize long-range correlations in genomic sequences.

\textbf{Density and clustering} (KDE, ScoreDBSCAN).
KDE applies Gaussian kernel density estimation to the score-weighted position profile and thresholds by percentile. ScoreDBSCAN maps positions to a 2D (coordinate, score) space and applies DBSCAN clustering with rescaled axes.
Density-based clustering was originally proposed by Ester et al.~\cite{ester1996dbscan} and has been adapted for mutation hotspot detection in cancer genomics; notably, Tokheim et al.~\cite{tokheim2016hotmaps} applied density-based clustering to 3D protein structures for exome-scale hotspot discovery.

\textbf{Statistical and model-based} (SlidingTest, LocalContrast, GradPeak, Bayes, SWAN).
These methods apply statistical tests to identify positions or windows significantly deviating from the background: SlidingTest and SWAN use sliding-window $z$-tests against the global mean, LocalContrast computes local $z$-scores, GradPeak detects prominence-filtered peaks on smoothed profiles, and Bayes computes posterior hotspot probabilities from a prior and normalized scores.
Sliding window approaches are a standard technique in population genetics and genomic analysis; Tajima~\cite{tajima1989statistical} introduced sliding window computation of the D statistic for detecting departures from neutrality, and similar windowed scanning strategies have been widely adopted for identifying regions under selection in viral genomes.

\textbf{Legacy} (MutClust-Orig).
MutClust-Orig applies the original MutClust logic (CCM seed + adaptive DBSCAN) to MSA-based score profiles, serving as a reference baseline representing the prior state-of-the-art in domain-specific hotspot detection~\cite{youn2025mutclust}.
Although MutClust uses DBSCAN clustering internally, it is separated from the Density/Clustering category because it constitutes a virus-specific pipeline (H-score computation $\to$ CCM seed selection $\to$ adaptive $\varepsilon$ clustering) rather than a generic density-based algorithm applied to arbitrary score profiles.
As the only existing virus-specific hotspot detection tool and the precursor study that motivated MutBench, it serves as the domain-specific reference baseline.
The detailed description of MutClust and the five Stage~1 variants is presented in Section~\ref{subsec:variants}.

%% ============================================================
\subsubsection{Experimental Design Considerations}
\label{sec:design_considerations}

The 11 pathogens in this benchmark differ in several structural characteristics that are not controlled across pathogens (Table~\ref{tab:design_characteristics}).

\begin{table}[htbp]
\centering
\caption[Structural characteristics of the 11 benchmark pathogens]{Structural characteristics of the 11 benchmark pathogens. AA lengths are post-truncation (Rabies G 491, EV-A71 VP1 261). Positive rates use operational alignment-QC Layer~A counts (HIV-1: 26/450 = 5.8\%); literature-curated counts (HIV-1: 23) in Table~\ref{tab:gt_positions} are the canonical narrative reference.}
\label{tab:design_characteristics}
\small
\begin{tabular}{lrrrr}
\toprule
\textbf{Pathogen} & \textbf{Unique seqs} & \textbf{AA length} & \textbf{Time span} & \textbf{Positive rate} \\
\midrule
SARS-CoV-2 & 753 & 1,259 & $\sim$5 yr & 2.0\% \\
H3N2 & 2,562 & 543 & $\sim$50 yr & 4.6\% \\
Norovirus & 791 & 536 & $\sim$20 yr & 2.8\% \\
HIV-1 & 2,256 & 450 & $\sim$30 yr & 5.8\% \\
Dengue & 796 & 490 & $\sim$20 yr & 4.9\% \\
RSV & 4,779 & 550 & $\sim$10 yr & 3.6\% \\
Influenza B & 5,019 & 560 & $\sim$10 yr & 3.0\% \\
MERS & 530 & 1,330 & $\sim$5 yr & 0.7\% \\
HCV & 1,067 & 340 & $\sim$30 yr & 15.9\% \\
Rabies & 2,038 & 491 & $\sim$50 yr & 7.9\% \\
EV-A71 & 662 & 261 & $\sim$20 yr & 10.3\% \\
\bottomrule
\end{tabular}
\end{table}

These differences---unique sequence count (530--5,019), protein length (261--1,330 AA), temporal coverage ($\sim$2--50 years), and Layer~A positive rate (0.7--15.9\%)---reflect the real-world heterogeneity of available viral sequence data rather than a controlled experimental design. The ``pathogen'' factor in the ANOVA is therefore a composite variable that conflates biological differences with data availability differences.

Three considerations mitigate this limitation:
(1)~subsampling robustness analysis (Section~\ref{sec:robustness_validation}) shows that MCC values stabilize at approximately 1,000 sequences and vary by less than $\pm 0.002$ at 25--100\% subsampling, indicating that sample size is not a dominant driver of performance differences;
(2)~the key finding is the \emph{scoring $\times$ pathogen interaction} ($\omega^2 = 0.296$), which operates within each pathogen where sample size and positive rate are held constant;
and (3)~equalizing sample sizes via subsampling to the minimum ($n = 530$, MERS) would sacrifice statistical power for pathogens with richer data without eliminating other confounds (temporal coverage, protein length).
A fully controlled design---equal-size, equal-quality MSAs spanning identical time ranges---is infeasible for real viral data but remains a useful benchmark for future synthetic studies.

%% ============================================================
\subsection{Evaluation and Statistical Design}
\label{sec:eval_stat_design}

\subsubsection{Evaluation Metrics}
\label{subsec:9pathogen_metrics}

The following evaluation metrics are used in the 11-pathogen benchmark.

\begin{itemize}
\item \textbf{MCC (Matthews Correlation Coefficient)}: The primary evaluation metric (defined formally in Equation~\ref{eq:mcc} below).
The confusion matrix is constructed using only Layer A (positive selection positions) as positive labels: TP = detected positions belonging to Layer A, FP = detected positions not belonging to Layer A, FN = undetected positions belonging to Layer A, TN = undetected positions not belonging to Layer A.
Layer B (constrained regions) falls within the standard MCC true-negative pool (Layer~B positions are non-Layer-A and therefore counted under TN if undetected, FP if detected); it is additionally tracked via a separate constrained false positive rate (constrained FPR) restricted to Layer~B. The majority of genomic positions belonging to neither Layer A nor B are similarly classified as FP or TN depending on detection status.
\item \textbf{F1}: The harmonic mean of precision and recall.
\item \textbf{Precision}: The proportion of true positives among detected positions.
\item \textbf{Recall}: The proportion of true positives detected among all true positives.
\item \textbf{Enrichment ratio}: The ratio of positive density within detected regions to overall positive density.
All enrichment tests used one-sided Fisher's exact test (alternative: greater), testing the directional hypothesis that hotspot positions are enriched for vaccine escape mutations relative to non-hotspot positions.
\item \textbf{Constrained FPR}: The false positive rate in Layer B constrained regions.
\item \textbf{DMS F1}: F1 score against Layer C DMS ground truth (computed only for 6 pathogens with DMS data).
\end{itemize}

Stage~1 uses the hotspot-score composite metric (defined in Eq.~\ref{eq:hotspot_score}, Section~\ref{subsec:hotspot_score}), which balances detection accuracy with reproducibility.
Stage~2 uses MCC as the primary metric instead, for three reasons:
(1)~the stability component of hotspot-score measures bootstrap reproducibility under position-level subsampling, but the subsampling behavior differs fundamentally across pathogens with different alignment lengths (340--1,330 AA) and MSA sizes (530--5,019 unique sequences), making cross-pathogen stability values non-comparable;
(2)~MCC utilizes all four cells of the confusion matrix (TP, TN, FP, FN), providing a balanced evaluation under the extreme class imbalance inherent in hotspot detection (positive rates 0.7\%--15.9\% across the 11 pathogens);
and (3)~separating the evaluation metric between stages allows Stage~1 to assess detection \textit{reliability} (via stability) and Stage~2 to assess detection \textit{accuracy} (via MCC), each optimized for its analytical purpose.
MCC is formally defined as:
\begin{equation}
\text{MCC} = \frac{TP \times TN - FP \times FN}{\sqrt{(TP+FP)(TP+FN)(TN+FP)(TN+FN)}}
\label{eq:mcc}
\end{equation}
where $TP$, $TN$, $FP$, and $FN$ denote true positives, true negatives, false positives, and false negatives, respectively.
MCC ranges from $-1$ (complete disagreement) through $0$ (random prediction) to $+1$ (perfect prediction), and unlike F1, it utilizes all four cells of the confusion matrix, making it robust to class imbalance.
When any factor in the denominator is zero (e.g., no positive predictions), MCC is conventionally defined as 0.
MCC was selected as the primary metric because in hotspot detection, where positive positions (convergent evolution sites) constitute a minority of the genome, it provides a more balanced evaluation than F1 for imbalanced classification problems. Chicco and Jurman~\cite{chicco2020mcc} systematically demonstrated that MCC provides more informative evaluation than F1 and accuracy in imbalanced binary classification.
The stability component used in Stage 1 is excluded from Stage 2 because it is an internal consistency metric independent of external ground truth, and cross-pathogen comparability of Bernoulli masking effects has not been validated given the wide variation in sequence length and MSA size across the 11 pathogens.

\subsubsection{Three-way ANOVA ($\omega^2$)}
\label{subsec:9pathogen_statistics}

A three-way ANOVA is performed with scoring type (20 levels), detection family (14 levels), and pathogen (11 levels) as factors.
The ANOVA used the 8{,}580 evaluation rows and includes all three main effects (scoring, family, pathogen) and all three two-way interactions (scoring $\times$ family, scoring $\times$ pathogen, family $\times$ pathogen). Detector-variant differences within each scoring--family--pathogen cell are retained in the residual term rather than modeled as a full scoring $\times$ family $\times$ pathogen interaction; this choice treats detector variants as within-family implementation noise and keeps the load-bearing inference on the pre-specified two-way interaction terms. The resulting residual degrees of freedom are 7{,}970, matching the archived \texttt{stage3\_statistics.csv}. The effective independent unit is nevertheless closer to the $20 \times 14 \times 11 = 3{,}080$ scoring--family--pathogen cells than to 8{,}580 fully independent observations, so Chapter~\ref{ch:results} reports pathogen-cluster bootstrap intervals for the headline interaction.
Omega-squared ($\omega^2$, Equation~\ref{eq:omega_squared}; non-partial form with $SS_{\text{total}} + MS_{\text{error}}$ as the denominator) was computed as the bias-corrected estimate of variance explained. Type~II sum of squares was used because it does not depend on the order of factor entry and provides unbiased estimates when interactions are present. Under Cohen's benchmarks for behavioral sciences~\cite{cohen1988statistical}, $\omega^2 > 0.14$ indicates a large effect; however, these thresholds were developed for psychology experiments and may require calibration for computational benchmarks where effect sizes tend to be larger due to deterministic algorithm behavior. The large evaluation count supports numerical stability of the fixed-effect fit, but inference is interpreted through the pathogen-cluster bootstrap because detector variants within a family are not fully independent.
Levene's test indicated heterogeneous variances across groups ($p < 0.05$); however, with $N = 8{,}580$ observations and a balanced design, ANOVA is robust to moderate violations of homoscedasticity.
As a safeguard, the ANOVA was accompanied by nonparametric checks.
The Kruskal--Wallis test applied to the scoring main effect yielded $H = 277.9$, $p < 10^{-47}$, consistent with the parametric main-effect result but not by itself testing interactions. An additional Aligned Rank Transform (ART) two-way ANOVA on the scoring~$\times$~pathogen interaction~\cite{wobbrock2011aligned} yielded a rank-scale interaction $\omega^2 = 0.277$ with $F_{190, 2860} = 7.33$ and $p = 8.6 \times 10^{-141}$, providing distribution-free support for the specific interaction that carries the main conclusion.
No family-wise multiple testing correction (e.g., Bonferroni, Holm) was applied to the ANOVA results because the analysis is exploratory rather than confirmatory: the goal is to decompose variance sources, not to test individual pairwise comparisons. The Friedman test and LOPO cross-validation are reported as additional lenses on the same dataset, not as independent replications. For the vaccine escape Fisher's exact tests, the per-pathogen family used in Chapter~\ref{ch:results} is the $20\times39 = 780$ scoring--detection combinations searched per pathogen (not 3), so Bonferroni-adjusted $p$-values are reported as $p_{\text{adj}} = \min(1, 780\,p)$; cross-pathogen denominators (joint $8{,}580$) are reported in the Table~\ref{tab:vaccine_escape_stage3} caption and do not change the three-tier tier assignment.
Following Cohen's effect size benchmarks~\cite{cohen1988statistical}, $\omega^2$ values are interpreted as small (0.01), medium (0.06), and large ($\geq 0.14$):
\begin{equation}
\omega^2 = \frac{SS_{\text{effect}} - df_{\text{effect}} \times MS_{\text{error}}}{SS_{\text{total}} + MS_{\text{error}}}
\label{eq:omega_squared}
\end{equation}
where $SS_{\text{effect}}$ is the sum of squares for the factor of interest, $df_{\text{effect}}$ its degrees of freedom, $MS_{\text{error}}$ the mean square of the residual error, and $SS_{\text{total}}$ the total sum of squares.
The $\omega^2$ values of main effects and two-way interactions are compared to identify the primary sources of performance variation.

\subsubsection{LOPO Cross-Validation (Leave-One-Pathogen-Out)}

The best-performing combination (scoring $\times$ detection) is selected from $n-1$ of 11 pathogens, applied to the held-out pathogen, and generalization performance is evaluated.
The best-performing combination is defined as the (scoring, family) pair with the highest mean MCC across the $n-1$ training pathogens.
Through 11 iterations, the concordance rate between training-best and test-best (oracle) combinations and the performance gap are computed.
Under random selection from 280 possible combinations ($20 \times 14$), the expected concordance rate is approximately 0.36\%.

\subsubsection{Friedman Test}

A nonparametric repeated-measures test using 11 pathogens as blocks, the top $k = 20$ scoring $\times$ family combinations were compared to test whether statistically significant performance differences exist among them.
Kendall's coefficient of concordance ($W$) is reported alongside $\chi^2$ to quantify ranking agreement across pathogens: $W = 1$ indicates perfect agreement, while $W = 0$ indicates no agreement.
Critical values and interpretation follow Dem\v{s}ar~\cite{demsar2006statistical}.

\subsubsection{Pathogen-level BCa Bootstrap Confidence Intervals}

Bias-corrected and accelerated (BCa) bootstrap confidence intervals (10,000 resamples) are computed for MCC values across 11 pathogens to quantify the performance variability of the best-performing combination across pathogens. With $n = 11$, the actual coverage of BCa intervals may be somewhat below the nominal 95\% level; these intervals should therefore be interpreted as approximate confidence intervals.
A separate combination-level cluster bootstrap (1{,}000 resamples, clustered by pathogen) is reported in Chapter~\ref{ch:results} to bound the spread of per-pathogen best-combination MCC under resampling of pathogens; the lower resample count reflects the much larger compute footprint of running the full $20\times39$ scoring--detection grid per resample, and is sufficient for stable percentile-CI estimation at $n_{\text{path}}=11$.

\subsubsection{Subsample Robustness Analysis}

Random subsamples at 25\%, 50\%, 75\%, and 100\% of MSA sequences are drawn, and the benchmark is repeated to verify MCC stability as a function of sample size.

\subsubsection{Feature Ablation Analysis}

To identify which information types contribute most to multi-source integration performance, a leave-one-feature-out ablation study is performed.
Starting from the 10-feature EqualWeight ensemble, each of the 10 scoring types is removed in turn, and the resulting 9-feature ensemble's mean MCC across 11 pathogens is compared to the full 10-feature baseline.
The change in mean MCC ($\Delta$MCC = MCC$_{\text{leave-one-out}}$ $-$ MCC$_{\text{full}}$) quantifies each feature's marginal contribution: negative $\Delta$ indicates that removing the feature \emph{lowers} ensemble performance (i.e., the feature contributes positively to the ensemble), while positive $\Delta$ indicates that the feature introduces net noise under uniform weighting and the ensemble is better off without it.
Additionally, nested feature subsets of sizes 2, 3, \ldots, 9 are evaluated (selecting features by descending individual AUC) to identify the minimum feature set that captures a target proportion (e.g., 98\%) of full-ensemble performance.

\textbf{Weight Sensitivity Analysis.}
\label{subsec:weight_sensitivity_methods}
To verify the robustness of assigning equal weights ($1/3$) to the three components of the hotspot-score, a generalized weighted hotspot-score is defined:
\begin{equation}
\text{hotspot-score}_w = w_r \cdot \text{recall} + w_p \cdot \text{precision} + w_s \cdot \text{stability}
\label{eq:weighted_hotspot_score}
\end{equation}
where $w_r + w_p + w_s = 1$ and each weight lies in the range $[0.05, 0.95]$.
Rankings of the five methods are computed across a total of 171 weight combinations to analyze the proportion of weight space in which MutClust-Hybrid maintains first place (rank robustness).
Additionally, rankings are separately presented for seven application scenarios, including early warning (recall-emphasis), vaccine design (precision-emphasis), and monitoring (stability-emphasis).

\textbf{Statistical Validation Methods.}
\label{subsec:statistical_methods}
The following three methods are applied to assess the statistical robustness of the benchmark results.

\textit{Bootstrap confidence intervals}: Region-level bootstrap is performed using the seven functional regions as units.
In each iteration, seven regions are sampled with replacement to construct a bootstrap ground truth, and the detection results from the full data are re-evaluated against this bootstrap ground truth.
This design captures the uncertainty in ground truth definition---variation in evaluation depending on which functional regions are included---and avoids the effective size reduction bias that occurs with individual position resampling.

\textit{Permutation test}: Two permutation approaches are employed.
(1)~Random label shuffling: Through 1,000 random permutations of ground truth labels, it is tested whether the observed metrics significantly exceed random expectation.
(2)~Circular permutation: The H-score vector is rotated by a random offset $k$ ($H_i \rightarrow H_{(i+k) \bmod N}$), hotspots are re-detected from the rotated scores, and evaluated against the original ground truth.
This approach preserves the spatial autocorrelation structure (cluster pattern) of H-scores while randomizing only the relationship with genomic positions, generating a more conservative null distribution than random label shuffling.
The p-value is calculated as $count / n_{\text{perm}}$ (the proportion of null distribution values that equal or exceed the observed value).

\textit{Pairwise comparison}: A paired bootstrap design that evaluates two methods simultaneously on the same bootstrap resample, testing the statistical significance of hotspot-score differences between methods.

\subsection{Stage~3: Multi-Source Information Integration Methods}
\label{subsec:stage3_methods}

Stage~3 layers an information-integration analysis on top of the Stage~2 evaluation matrix. The same 20 scoring types are decomposed into 10 underlying positional features (homoplasy, entropy, frequency, rare-variant frequency, dN/dS proxy, pLDDT, SASA, Grantham distance, ESM-2 masked-marginal score, semantic change), and three ablation analyses are performed on an EqualWeight ensemble of those 10 features:
(i)~\emph{leave-one-feature-out}: each feature is removed in turn and the change in mean MCC ($\Delta$MCC) is reported;
(ii)~\emph{group ablation}: features are partitioned into Frequency-only (3), Freq~+~Phylogeny (4), Freq~+~Structure (5), Freq~+~DL (5), and All-10, and group-wise mean MCC is compared;
(iii)~\emph{cumulative addition}: features are added in descending order of single-feature AUC and the cumulative mean MCC is computed at $k = 1, 2, \ldots, 10$ to identify the minimum subset that recovers a target proportion of the All-10 performance.
Per-pathogen feature AUC is computed via threshold sweeping (no model training), and a complementary multivariate Random Forest (\texttt{n\_estimators}~=~200, \texttt{max\_depth}~=~5, \texttt{class\_weight}~=~\texttt{balanced}, \texttt{random\_state}~=~42; stratified 5-fold cross-validation per pathogen; cf.\ Section~\ref{subsec:rf_importance}) reports relative feature importances for triangulation.
The Stage~3 evaluation set is the same 11 pathogens used in Stage~2, optionally extended to 12 pathogens by adding Zika (E protein, 505~AA) when the analysis benefits from increased per-feature statistical power; the 11-vs-12 distinction is reported transparently in every Stage~3 table caption (Section~\ref{subsec:feature_ablation}).
A vaccine-escape cross-validation analysis is then performed on the 3 of 11 pathogens with curated antibody-escape annotations (H3N2, SARS-CoV-2, HIV-1; 2{,}340 enrichment evaluations) and audited for circularity against Layer~A (Section~\ref{sec:cross_validation_escape}).

\subsection{Code, Data, and Reproducibility}
\label{subsec:code_data_availability}
The full source code, configuration files, scoring scripts, detection-family implementations, ground-truth tables, statistical-analysis notebooks, and the 8{,}580-evaluation result archive are released under an open-source license (see Chapter~\ref{ch:conclusion}, Section~\ref{sec:code_availability} for the canonical repository URL, commit hash, and Zenodo DOI). The pinned random seed (\texttt{seed} = 42) governs all bootstrap, subsampling, and permutation analyses; library versions (Python~3.10, NumPy~1.24, SciPy~1.11, IQ-TREE~2, HyPhy~v2.5, ESM-2 \texttt{esm2\_t33\_650M\_UR50D}) are pinned in the released environment file.

\bigskip
\noindent\textbf{Chapter summary.}
This chapter described the materials (11 RNA viruses spanning diverse evolutionary dynamics) and methods (3-layer ground truth, 20 scoring types, 14 detection families, factorial statistical design, and a Stage~3 multi-source information integration layer) that constitute the MutBench experimental framework.
The following chapter applies this framework to report experimental results: Stage~1 provides a controlled SARS-CoV-2 sanity check, Stage~2 extends the evaluation to all 11 pathogens to identify the dominant factors influencing hotspot detection performance, and Stage~3 decomposes the resulting performance by information type via per-feature AUC, ablation, and Random Forest analyses.
